% =============================================================================
%  PRISMA 2020 — Checklist de 27 Itens
%  Versão em Português Brasileiro (PT-BR)
%  Compilar com: pdflatex prisma2020-checklist.tex
% =============================================================================
\documentclass[11pt]{article}

% ─── Pacotes ─────────────────────────────────────────────────────────────────
\usepackage[utf8]{inputenc}
\usepackage[T1]{fontenc}
\usepackage[brazil]{babel}
\usepackage{geometry}
\geometry{a4paper, margin=1.8cm}
\usepackage{xcolor}
\usepackage{array}
\usepackage{longtable}
\usepackage{booktabs}
\usepackage{hyperref}
\hypersetup{
  colorlinks=true,
  linkcolor=blue!60!black,
  urlcolor=blue!60!black,
  citecolor=blue!60!black
}

% ─── Cores ───────────────────────────────────────────────────────────────────
\definecolor{corSecao}{RGB}{41, 82, 163}
\definecolor{corSub}{RGB}{70, 130, 180}

% ─── Estilo da Tabela ───────────────────────────────────────────────────────
\newcolumntype{C}[1]{>{\centering\arraybackslash}p{#1}}
\newcolumntype{L}[1]{>{\raggedright\arraybackslash}p{#1}}

% Comando para seções da checklist
\newcommand{\secao}[1]{%
  \rowcolor{corSecao!12}%
  \multicolumn{4}{L{\dimexpr\textwidth-2\tabcolsep\relax}}{%
    \textbf{\textcolor{corSecao}{#1}}}%
  \tabularnewline%
}

% Comando para subseção (tipo "Reporting Item")
\newcommand{\subsecao}[1]{%
  \rowcolor{corSub!8}%
  \multicolumn{4}{L{\dimexpr\textwidth-2\tabcolsep\relax}}{%
    \textit{#1}}%
  \tabularnewline%
}

% ─── Cabeçalho Repetido ─────────────────────────────────────────────────────
\newenvironment{prismachecklist}{%
\begin{longtable}{
  >{\footnotesize\bfseries}C{1.2cm}
  >{\footnotesize}L{2.8cm}
  >{\footnotesize}L{9cm}
  >{\footnotesize\centering}C{2cm}
}
  \toprule
  \rowcolor{corSecao!20}
  \textbf{\#} &
  \textbf{Seção / Tópico} &
  \textbf{Descrição do item} &
  \textbf{P.\ rel.} \\
  \midrule
  \endfirsthead
  \toprule
  \rowcolor{corSecao!20}
  \textbf{\#} &
  \textbf{Seção / Tópico} &
  \textbf{Descrição do item} &
  \textbf{P.\ rel.} \\
  \midrule
  \endhead
  \midrule
  \multicolumn{4}{r}{\footnotesize\itshape Continua na próxima página...} \\
  \endfoot
  \bottomrule
  \endlastfoot
}{%
  \end{longtable}
}

% =============================================================================
\begin{document}
\thispagestyle{empty}
\begin{center}
  {\Large\bfseries PRISMA 2020 --- Checklist de 27 Itens}\\[4pt]
  {\small\itshape Preferred Reporting Items for Systematic Reviews and
           Meta-Analyses}\\[4pt]
  {\small Preencha a coluna ``P.\ rel.'' (página relatada) conforme a
         localização no manuscrito.}
\end{center}

\vspace{0.4cm}

\begin{prismachecklist}

% ══════════════════════════════════════════════════════════════════════════════
%  TÍTULO
% ══════════════════════════════════════════════════════════════════════════════
\secao{TÍTULO --- \textit{TITLE}}

1 & Título &
  \textbf{Título.} Identifique o artigo como uma revisão sistemática,
  meta-análise, ou ambos. &
  \\
  \addlinespace

% ══════════════════════════════════════════════════════════════════════════════
%  RESUMO
% ══════════════════════════════════════════════════════════════════════════════
\secao{RESUMO --- \textit{ABSTRACT}}

2 & Resumo estruturado &
  \textbf{Resumo estruturado.} Forneça um resumo estruturado que inclua:
  histórico; objetivos; critérios de elegibilidade; fontes de informação;
  risco de viés; resultados; discussão; registro e financiamento. &
  \\
  \addlinespace

% ══════════════════════════════════════════════════════════════════════════════
%  INTRODUÇÃO
% ══════════════════════════════════════════════════════════════════════════════
\secao{INTRODUÇÃO --- \textit{INTRODUCTION}}

3 & Justificativa &
  \textbf{Justificativa.} Descreva a justificativa para a revisão no
  contexto do conhecimento existente. &
  \\

4 & Objetivos &
  \textbf{Objetivos.} Forneça uma declaração explícita dos objetivos ou
  perguntas que a revisão aborda. &
  \\
  \addlinespace

% ══════════════════════════════════════════════════════════════════════════════
%  MÉTODOS
% ══════════════════════════════════════════════════════════════════════════════
\secao{MÉTODOS --- \textit{METHODS}}

5 & Critérios de elegibilidade &
  \textbf{Critérios de elegibilidade.} Especifique os critérios de inclusão
  e exclusão para a revisão e como os estudos foram agrupados para as
  sínteses. &
  \\

6 & Fontes de informação &
  \textbf{Fontes de informação.} Especifique todas as bases de dados,
  registros, websites, organizações, listas de referência e outros locais
  onde a busca foi realizada ou atualizada. Informe a data da última busca
  em cada fonte. &
  \\

7 & Estratégia de busca &
  \textbf{Estratégia de busca.} Apresente as estratégias de busca completas
  para todas as bases de dados e registros, incluindo filtros e limites
  utilizados. &
  \\

8 & Processo de seleção &
  \textbf{Processo de seleção.} Especifique os métodos usados para decidir
  se um estudo atende aos critérios de inclusão, incluindo quantos
  revisores analisaram cada registro e cada relato, se o trabalho foi
  realizado de forma independente e, se aplicável, os detalhes das
  ferramentas de automação utilizadas. &
  \\

9 & Processo de coleta de dados &
  \textbf{Processo de coleta de dados.} Descreva os métodos usados para
  extrair dados dos relatos, incluindo quantos revisores extraíram dados
  de cada relato, se trabalharam de forma independente, os processos
  para obter ou confirmar dados com os investigadores do estudo e, se
  aplicável, detalhes das ferramentas de automação. &
  \\

10 & Itens de dados &
  \textbf{Itens de dados.} Liste e defina todos os desfechos para os quais
  os dados foram extraídos. Liste e defina todas as outras variáveis para
  as quais os dados foram extraídos (características do estudo, fontes de
  financiamento). Descreva quaisquer suposições ou simplificações feitas. &
  \\

11 & Risco de viés dos estudos &
  \textbf{Risco de viés dos estudos.} Especifique os métodos usados para
  avaliar o risco de viés nos estudos incluídos, incluindo detalhes das
  ferramentas utilizadas, quantos revisores avaliaram cada estudo e se
  trabalharam de forma independente. &
  \\

12 & Medidas de efeito &
  \textbf{Medidas de efeito.} Especifique para cada desfecho a medida de
  efeito utilizada (risco relativo, diferença média, etc.). &
  \\

13 & Métodos de síntese &
  \textbf{Métodos de síntese.} Descreva o processo de síntese e análise
  dos resultados: elegibilidade para síntese, métodos de preparação dos
  dados, métodos de análise estatística, métodos de exploração de
  heterogeneidade e análises de sensibilidade. &
  \\

14 & Viés de relato &
  \textbf{Avaliação do viés de relato.} Descreva os métodos usados para
  avaliar o risco de viés devido a resultados não relatados (viés de
  relato seletivo). &
  \\

15 & Avaliação da certeza &
  \textbf{Avaliação da certeza da evidência.} Descreva os métodos usados
  para avaliar a certeza (confiança) no conjunto de evidências para cada
  desfecho (ex.: GRADE). &
  \\
  \addlinespace

% ══════════════════════════════════════════════════════════════════════════════
%  RESULTADOS
% ══════════════════════════════════════════════════════════════════════════════
\secao{RESULTADOS --- \textit{RESULTS}}

16 & Seleção dos estudos &
  \textbf{Seleção dos estudos.} Descreva os resultados do processo de
  busca e seleção, desde o número de registros identificados até o número
  de estudos incluídos na revisão, idealmente usando um diagrama de fluxo. &
  \\

17 & Características dos estudos &
  \textbf{Características dos estudos.} Cite cada estudo incluído e
  apresente suas características (tamanho da amostra, duração do
  acompanhamento, tipo de intervenção, etc.). &
  \\

18 & Risco de viés nos estudos &
  \textbf{Risco de viés nos estudos.} Apresente as avaliações de risco de
  viés para cada estudo incluído. &
  \\

19 & Resultados individuais &
  \textbf{Resultados individuais.} Para cada desfecho, apresente, de forma
  ideal: estimativas de efeito e intervalos de confiança para cada estudo,
  idealmente com gráficos de floresta ou mapas de resultados. &
  \\

20 & Resultados da síntese &
  \textbf{Resultados da síntese.} Para cada desfecho, apresente os
  resultados de todas as sínteses estatísticas realizadas. Se
  meta-análise foi feita: apresente para cada uma a estimativa de efeito
  agregado e seu intervalo de confiança/previsão, e medidas de
  heterogeneidade. &
  \\

21 & Viés de relato &
  \textbf{Viés de relato.} Apresente as avaliações do risco de viés de
  relato seletivo. &
  \\

22 & Certeza da evidência &
  \textbf{Certeza da evidência.} Apresente as avaliações da certeza
  (confiança) no conjunto de evidências para cada desfecho. &
  \\
  \addlinespace

% ══════════════════════════════════════════════════════════════════════════════
%  DISCUSSÃO
% ══════════════════════════════════════════════════════════════════════════════
\secao{DISCUSSÃO --- \textit{DISCUSSION}}

23 & Discussão &
  \textbf{Discussão.} Forneça uma interpretação geral dos resultados no
  contexto de outras evidências e implicações para prática, pesquisa e
  políticas. Discuta limitações da evidência incluída e do processo de
  revisão. &
  \\
  \addlinespace

% ══════════════════════════════════════════════════════════════════════════════
%  OUTRAS INFORMAÇÕES
% ══════════════════════════════════════════════════════════════════════════════
\secao{OUTRAS INFORMAÇÕES --- \textit{OTHER INFORMATION}}

24 & Registro e protocolo &
  \textbf{Registro e protocolo.} Informe onde a revisão foi registrada e
  o identificador do registro. Indique se um protocolo foi preparado e
  onde pode ser acessado. Descreva quaisquer alterações ao protocolo. &
  \\

25 & Suporte &
  \textbf{Fonte de financiamento.} Descreva as fontes de suporte financeiro
  e não financeiro para a revisão, o papel dos financiadores e quaisquer
  conflitos de interesse. &
  \\

26 & Conflitos de interesse &
  \textbf{Conflitos de interesse.} Declare quaisquer conflitos de interesse
  dos autores da revisão. &
  \\

27 & Disponibilidade de dados &
  \textbf{Disponibilidade de dados, código e outros materiais.} Indique
  quais dos seguintes itens estão publicamente disponíveis e onde podem
  ser encontrados: dados de extração, código de análise, outros materiais
  usados na revisão. &
  \\

\end{prismachecklist}

\vspace{0.5cm}

\begin{minipage}{0.85\textwidth}
\small
\textbf{Citação:} Page MJ, McKenzie JE, Bossuyt PM, et al. \textit{The PRISMA
2020 statement: an updated guideline for reporting systematic reviews.}
BMJ 2021;372:n71. \href{https://doi.org/10.1136/bmj.n71}{doi:
10.1136/bmj.n71}\\[4pt]
\textbf{Licença:} CC-BY 4.0. Este checklist é uma adaptação para o
português brasileiro do original PRISMA 2020.
\end{minipage}

\end{document}
